\section{Regularization and angular control}\label{sec:energy}

We now combine quadratic mass growth with regularized current
pairings to prove that \(Q\in L^2\) and to control the variation
of nearby complex-line directions. These are the remaining
quantitative inputs to the density-splitting theorem.

Let \(\Delta=dd^*+d^*d\) be the Hodge Laplacian, and set
\[
S_r=(1+r^2\Delta)^{-3},\qquad 0<r<r_0.
\]
The exponent is fixed. By Proposition~\ref{prop:separation},
the regularizations of \(P,Q,T\) belong to \(L^2\):
for fixed \(r\), \(S_r\) has order \(-6\) and maps
\(H^{-3}\) into \(H^3\).
The same is true for a positive current \(P'\) represented
by any \(\lambda'\le\lambda\).
Write \(\mu'\) for the projection of \(\lambda'\).

\begin{lemma}\label{lem:zero-pairing}
If \(T'\) is a closed current with \(S_rT'\in L^2\), then
\begin{equation}\label{eq:zero-pairing}
\ip{S_rT}{*S_rT'}_{L^2}=0.
\end{equation}
\end{lemma}

\begin{proof}
The operator \(S_r\) is self-adjoint, commutes with star in
middle degree, and commutes across degrees with \(d\) and
\(d^*\). It preserves smooth forms. Thus \(S_rT\) annihilates
every smooth closed form by \eqref{eq:annihilation}.
Closedness of \(T'\) as a current means \(d^*T'=0\), so
\(*S_rT'\) is a distributionally closed \(L^2\) form.
Its heat regularizations
are smooth closed forms and converge to it in \(L^2\).
The asserted pairing follows by continuity.
\end{proof}

\subsection{The kernel and its positive leading part}

For nearby \(x,y\), let \(\tau_{yx}\) denote metric parallel
transport along the short geodesic from \(y\) to \(x\).
\begin{lemma}\label{lem:kernel}
The continuous kernel \(K_r(x,y)\) of \(S_r^2\) can be written
\[
K_r(x,y)=p_r(x,y)\tau_{yx}+E_r(x,y),
\]
where \(p_r\ge0\) is supported in a fixed neighborhood of the
diagonal. For every fixed positive integer \(N\),
\begin{align}
0\le p_r(x,y)&\le C_Nr^{-4}(1+d_g(x,y)/r)^{-N},
\label{eq:p-upper}\\
p_r(x,y)&\ge cr^{-4}\qquad(d_g(x,y)<r),
\label{eq:p-lower}\\
|E_r(x,y)|&\le C_Nr^{-2}(1+d_g(x,y)/r)^{-N}.
\label{eq:error-kernel}
\end{align}
\end{lemma}

\begin{proof}
Spectral calculus gives
\begin{equation}\label{eq:gamma}
S_r^2=\frac1{\Gamma(6)}
\int_0^\infty e^{-s}s^5e^{-r^2s\Delta}\dd s.
\end{equation}
The Hodge Laplacian is the Levi-Civita connection Laplacian
on forms plus a zeroth-order curvature endomorphism.
The leading term of its heat kernel, localized by a
nonnegative diagonal cutoff, is
\[
\chi(x,y)(4\pi t)^{-2}e^{-d_g(x,y)^2/(4t)}
a(x,y)\tau_{yx},
\]
where \(a\) is a smooth positive scalar volume correction
and \(\chi=1\) sufficiently near the diagonal.
For the leading coefficient, see the author version
\cite[Section~3, equation~(3.4)]{Ludewig2018}; for
uniform short-time asymptotics and the global Gaussian bound,
see \cite[Theorems~1.1 and~3.5]{Ludewig2019}.

Fix a small \(t_0>0\). After subtracting the displayed term,
the short-time remainder has the bound
\begin{equation}\label{eq:heat-remainder}
Ct^{-1}e^{-c_0d_g(x,y)^2/t}+C_Mt^M,\qquad 0<t<t_0,
\end{equation}
for any prescribed integer \(M\), with constants depending
on \(M\). Indeed, take a sufficiently long heat parametrix:
the further localized coefficients gain at least one factor
of \(t\), and its final uniform remainder can be made
\(O(t^M)\). Off the diagonal the cutoff errors are of
arbitrarily high order.

Integrating the leading term for \(r^2s<t_0\) in
\eqref{eq:gamma} defines \(p_r\). For each fixed \(N\),
\[
e^{-c_0d^2/(r^2s)}
\le C_N(1+\sqrt{s})^N(1+d/r)^{-N}.
\]
The resulting \(s\)-integrals converge: their powers at zero
are \(s^3\) for the leading term and \(s^4\) for the Gaussian
remainder, and \(e^{-s}\) controls infinity. This gives
\eqref{eq:p-upper} and the Gaussian part of
\eqref{eq:error-kernel}. Restricting to \(1\le s\le2\)
gives \eqref{eq:p-lower} for small \(r\).

The smooth remainder in \eqref{eq:heat-remainder} integrates
to \(O(r^{2M})\). Since \(X\) has bounded diameter, choosing
\(2M\ge N-2\) absorbs it into \eqref{eq:error-kernel}.
For \(s\ge t_0/r^2\), the heat kernel is uniformly bounded,
and the tail of the gamma weight is exponentially small.
This proves all the bounds. Notice that increasing \(N\)
changes the parametrix length, not the fixed resolvent
power.
\end{proof}

For unit complex line bivectors at \(x,y\), respectively
\(\ell,m\), the algebra \eqref{eq:algebra} gives the exact
identity
\begin{equation}\label{eq:angular-identity}
\ip{\ell}{*\tau_{yx}m}
=\tfrac12\ip{F_x}{\tau_{yx}F_y}-\ip{\ell}{\tau_{yx}m}
=\tfrac12|\ell-\tau_{yx}m|^2
 -\tfrac14|F_x-\tau_{yx}F_y|^2.
\end{equation}
Here parallel transport is an isometry and commutes with star.
In particular the negative term is \(O(d_g(x,y)^2)\).

\begin{lemma}\label{lem:positive-pairing}
Uniformly for all \(\lambda'\le\lambda\),
\begin{equation}\label{eq:smoothed-mass}
\|S_rP'\|_2\le Cr^{-1},
\end{equation}
and, writing \(d=d_g(x,y)\),
\begin{align}
\ip{S_rP}{*S_rP'}_{L^2}
&\ge cr^{-4}\iint_{d<r}
|\ell-\tau_{yx}m|^2\dd\lambda(x,\ell)\dd\lambda'(y,m)
\nonumber\\
&\quad-Cr^{-2}\iint(1+d/r)^{-6}\dd\mu(x)\dd\mu'(y).
\label{eq:positive-pairing}
\end{align}
Moreover
\begin{equation}\label{eq:shell}
r^{-2}\int_X(1+d_g(x,y)/r)^{-6}\dd\mu(x)\le C.
\end{equation}
\end{lemma}

\begin{proof}
The inner ball and the dyadic shells in
Proposition~\ref{prop:growth} give \eqref{eq:shell}:
the shell at distance comparable to \(2^jr\) contributes
at most \(Cr^2\,2^{-4j}\) before division by \(r^2\).

Self-adjointness and commutation with star give
\[
\ip{S_rP}{*S_rP'}=
\iint\ip{\ell}{*K_r(x,y)m}
\dd\lambda(x,\ell)\dd\lambda'(y,m).
\]
The identity is justified directly by the continuous kernel,
or by approximation in \(H^{-3}\): the operator \(S_r^2\)
has order \(-12\) and maps \(H^{-3}\) into \(H^9\),
so the dual pairings are continuous. Using
\eqref{eq:angular-identity}, retain its nonnegative term
only on \(d<r\), where \eqref{eq:p-lower} applies.
By \eqref{eq:p-upper} with \(N\ge8\), the negative term
is bounded by
\(Cr^{-2}(1+d/r)^{-6}\); the kernel error has the same
bound. This proves \eqref{eq:positive-pairing}.
For \eqref{eq:smoothed-mass}, use the analogous identity
for \(\|S_rP'\|_2^2\), take absolute values of the kernel,
and apply \eqref{eq:shell} and \(\mu'\le\mu\).
\end{proof}

The kernel estimate and quadratic growth bound the negative error
in \eqref{eq:positive-pairing} uniformly. To extract an \(L^2\)
bound for \(Q\), we must also control the cross terms between
invariant and anti-invariant components. The next estimate
supplies this control.

\subsection{Mixing of invariant and anti-invariant components}

\begin{lemma}[Relative projection estimate]\label{lem:mixing}
Let \(\Pi\) be a fixed smooth orthogonal projection on
\(\Lambda^2T^*X\). If a distributional two-form \(W\)
satisfies \(\Pi W=0\) and \(S_rW\in L^2\), then
\begin{equation}\label{eq:mixing}
\|\Pi S_rW\|_2\le Cr\|S_rW\|_2.
\end{equation}
The constant may depend on \(\Pi\), but is independent of
\(r\) and \(W\).
\end{lemma}

\begin{proof}
Put \(U=S_rW\) and \(A_r=(1+r^2\Delta)^3\). The operators
\(S_r\) and \(A_r\) are inverse on distributions. Since
\(\Pi A_rU=\Pi W=0\), we have, distributionally,
\[
\Pi U=S_rA_r(\Pi U)=S_r[A_r,\Pi]U.
\]
On smooth forms the adjoint of the operator on the right
is \(-[A_r,\Pi]S_r\). Expanding the commutator gives terms
\(\binom3j r^{2j}[\Delta^j,\Pi]S_r\), \(1\le j\le3\).
The principal symbol of \(\Delta^j\) is scalar, so
\([\Delta^j,\Pi]\) has differential order at most \(2j-1\).
The elliptic Sobolev calculus
\cite[Propositions~6.7--6.8]{MelroseIML}, together with
spectral calculus for the Hodge Laplacian, gives
\[
\|S_r\|_{L^2\to H^{2j-1}}\le Cr^{-(2j-1)}.
\]
Each term of the adjoint therefore has norm \(O(r)\).
Taking its bounded Hilbert-space adjoint extends
\(S_r[A_r,\Pi]\) to \(L^2\), agreeing with its distributional
action. Applying this bound to \(U\) proves
\eqref{eq:mixing}.
\end{proof}

\begin{proposition}\label{prop:energy}
The correction \(Q\) belongs to \(L^2\), and
\begin{equation}\label{eq:angular}
\ang_r:=
\iint_{d_g(x,y)<r}|\ell-\tau_{yx}m|^2
\dd\lambda(x,\ell)\dd\lambda(y,m)\le Cr^4.
\end{equation}
For every positive current \(P'\) represented by
\(\lambda'\le\lambda\),
\begin{equation}\label{eq:cross-vanish}
\lim_{r\downarrow0}\ip{S_rP'}{*S_rQ}_{L^2}=0.
\end{equation}
\end{proposition}

\begin{proof}
For any anti-invariant distribution \(V\) with \(S_rV\in L^2\),
split both factors along \(R\) and \(1-R\), and use
Lemmas~\ref{lem:positive-pairing} and \ref{lem:mixing}:
\begin{equation}\label{eq:cross-bound}
|\ip{S_rP'}{S_rV}|
\le Cr\|S_rP'\|_2\|S_rV\|_2
\le C\|S_rV\|_2.
\end{equation}
Since \(Q\) is self-dual, Lemma~\ref{lem:zero-pairing} with
\(T'=T\) gives
\[
0=\ip{S_rP}{*S_rP}
+2\ip{S_rP}{S_rQ}+\|S_rQ\|_2^2.
\]
Equations~\eqref{eq:positive-pairing}, \eqref{eq:shell},
and \eqref{eq:cross-bound} imply
\[
cr^{-4}\ang_r+\|S_rQ\|_2^2
\le C+C\|S_rQ\|_2.
\]
Thus \(\|S_rQ\|_2\) is uniformly bounded and
\(\ang_r\le Cr^4\). A weakly convergent \(L^2\) subsequence,
together with distributional convergence \(S_rQ\to Q\),
proves \(Q\in L^2\).

For any anti-invariant \(V\in L^2\),
\eqref{eq:cross-bound} and the \(L^2\) contraction property
of \(S_r\) give a uniform bound by \(C\|V\|_2\).
If \(V\) is smooth, then it is self-dual and
\[
\ip{S_rP'}{*S_rV}=P'(S_r^2V)\longrightarrow P'(V)=0,
\]
because \(S_r^2V\to V\) smoothly. Smooth anti-invariant
forms are dense in the anti-invariant \(L^2\) subspace:
smooth first and then apply \(R\).
Approximation of \(Q\) now proves \eqref{eq:cross-vanish}.
\end{proof}
